\label{Section:result section}
In this section, we first present the results obtained for the \texttt{memcmps} code.\\
Next, we show the results for the \texttt{verify\_pin} code.\\
Finally, we discuss the findings related to the control-flow graph generated during the analysis.

%------------------------------------------------verify_pin-----------------------------------------------
\subsection{Verify pin}
This section presents the results for the \texttt{verifypin} test.

\begin{table*}[t]
\centering
\begin{tabular}{l|llll|llll|l}
Version    & \multicolumn{4}{c|}{Attacks} & \multicolumn{4}{c|}{Non-redundant attacks} & Countermeasures \\
           & 1F & 2F & 3F & 4F & 1F & 2F & 3F & 4F &                 \\
\hline
\hline
VP0        & 12 & 21  & 18  & 7    & 12 & 0  & 0   & 0   & $\emptyset$           \\
VP0 - rust & 8  & 20  & 20  & 10   & 8  & 0  & 0   & 0   & $\emptyset$           \\
\hline
VP1        & 12 & 21  & 18  & 7    & 12 & 0  & 0   & 0   & HB                    \\
VP1 - rust & 12 & 21  & 18  & 7    & 12 & 0  & 0   & 0   & HB                    \\
\hline
VP2        & 34 & 122 & 204 & 183  & 34 & 4  & 0   & 0   & HB+FTL                \\
VP2 - rust & 34 & 122 & 204 & 183  & 34 & 4  & 0   & 0   & HB+FTL                \\
\hline
VP3        & 35 & 122 & 204 & 183  & 35 & 4  & 0   & 0   & HB+FTL+INL            \\
VP3 - rust & 35 & 122 & 204 & 183  & 35 & 4  & 0   & 0   & HB+FTL+INL            \\
\hline
VP4        & 34 & 118 & 180 & 147  & 34 & 0  & 4   & 0   & FTL+INL+DPTC+PTCBK+LC \\
VP4 - rust & 34 & 118 & 180 & 147  & 34 & 0  & 4   & 0   & FTL+INL+DPTC+PTCBK+LC \\
\hline
VP5        & 0  & 80  & 644 & 2566 & 0  & 80 & 248 & 184 & HB+FTL+DPTC+DC        \\
VP5 - rust & 0  & 80  & 644 & 2566 & 0  & 80 & 248 & 184 & HB+FTL+DPTC+DC        \\
\hline
VP6        & 0  & 3   & 9   & 20   & 0  & 3  & 9   & 14  & HB+FTL+INL+DPTC+DT    \\
VP6 - rust & 0  & 3   & 9   & 20   & 0  & 3  & 9   & 20  & HB+FTL+INL+DPTC+DT    \\
\hline
VP7        & 1  & 2   & 8   & 13   & 1  & 2  & 8   & 13  & HB+FTL+INL+DPTC+DT+SC \\
VP7 - rust & 1  & 2   & 8   & 13   & 1  & 2  & 8   & 13  & HB+FTL+INL+DPTC+DT+SC
\end{tabular}
    \caption{Comparison of attacks between \texttt{verify\_pin} implementations \label{tbl:results-fissc-vp}}
\end{table*}

\noindent Table \ref{tbl:results-fissc-vp}: present the attacks identified by Lazart for the code C and Rust.
The column "Version" indicates the version of the program and the implementation language. 
The second to fifth columns correspond respectively to the number of attacks reported by Lazart for 1 fault to 4 faults.
Similarly, the next four columns present the number of \textit{non-redundant} attacks for 1 fault to 4 faults. 
The last column "Countermeasures" enumerates the protection used for each version (the protection are the same for version X, and it's Rust counterpart):
\begin{itemize}
        \item \textit{HB} : hardened boolean
        \item \textit{FTL} : fixed time loops
        \item \textit{INL} : inlined function
        \item \textit{DPTC} : try counter decremented first
        \item \textit{PTCBK} : try counter backup
        \item \textit{LC} : loop counter verification
        \item \textit{DC} : double call
        \item \textit{DT} : double test
        \item \textit{SC} : step counter
\end{itemize}

In the VP0 example, we observe 8 missing attacks when we compare to the C result. Since VP0 doesn't have hardened boolean, so it uses regular boolean, \texttt{rustc} reuse the branch condition in the \texttt{verify\_pin} function, through XOR operations, to compute its return value, reducing the number of observable injection points on that path. All the higher-order attacks are redundant with those 8 one fault attacks.

The \textit{redundancy analysis} regroups attacks that are variations of a simpler attack adding additional faults.

In the VP6 example, the non-redundant attacks count differs in 4-fault (while being the same without redundancy analysis) because \texttt{rustc} places some injection points earlier in the control flow. The IPs are therefore identified as the non-redundant ones, which changes the higher-order results.
%------------------------------------------------memcompare-----------------------------------------------
\subsection{memcmps}
\label{section:memcmps_result}
This section presents the results for the \texttt{memcmps} test.

\begin{table*}[t]
\centering
\begin{tabular}{l|llll|llll}
Version      & \multicolumn{4}{c|}{Attacks} & \multicolumn{4}{c}{Non-redundant attacks} \\
             & 1F & 2F & 3F & 4F & 1F & 2F & 3F & 4F \\
\hline
\hline
MCMPS1        & 16  & 32  & 28  & 8  & 16  & 0  & 0  & 0  \\
MCMPS1 - rust & 16  & 32  & 28  & 8  & 16  & 0  & 0  & 0  \\
\hline
MCMPS2        & 4  & 32  & 404  & 1400 & 4  & 32  & 322  & 194  \\
MCMPS2 - rust & 4  & 32  & 404  & 1400 & 4  & 32  & 322  & 194  \\
\hline
MCMPS3        & 4  & 32  & 216  & 2262   & 4  & 32  & 134  & 1468  \\
MCMPS3 - rust & 4  & 32  & 216  & 2262   & 4  & 32  & 134  & 1468  \\
\end{tabular}
\caption{Comparison of attacks between \texttt{memcmps} implementations \label{tbl:results-fissc-mcmp}}
\end{table*}

Table \ref{tbl:results-fissc-mcmp} indicates for each version and implementation language (first column "Version") the results obtained for 1 to 4 faults - as for Table \ref{tbl:results-fissc-vp}, the column "Attacks" shows attacks count while the "Non-redundant attacks" column shows the subset of non-redundant attacks.


Table \ref{tbl:results-fissc-mcmp} presents the attack found by Lazart for the code C and Rust.\\
Version 1 (V1) uses one mask and one call to \texttt{memcmp}.\\
Version 2 (V2) uses one mask and three calls to \texttt{memcmp}.\\
Version 3 (V3) uses two masks and four calls to \texttt{memcmp}.\\
Version 3 is also the one shown in Figure \ref{tbl:results-fissc-mcmp}.\\

As shown in Table \ref{tbl:results-fissc-mcmp}, the Rust implementation behaves exactly as the C implementation, showing the same result for 1 to 4 faults with a combination of TI and DL fault models.

%------------------------------------------------CFG-----------------------------------------------

       
\subsection{Control Flow Graph}
\label{subsection:result cfg}
We also observed some major differences in how Rust compiles compared to C.\\
The intermediate code generated by LLVM revealed that Rust optimises certain conditional blocks differently than C.\\

However, we noticed differences in the control-flow graphs generated by the two language versions of the code presented in this study.\\
Specifically, Rust transforms variables with two possible values—depending on the evaluation of a condition, into a \texttt{select} instruction rather than a branch instruction.
Similarly, Rust may use a \texttt{switch} instruction when a variable is compared to multiple values.